\chapter{提案手法}
\label{ch:app}
\quad

アノテーションの定義と，アノテーションの内部データ構造を説明する~\cite{flanagan2008ruby}．

\section{提案する手法 ・ アプローチ}
\label{sec:app_xx}

モゲという概念を用いる．モゲという概念を用いる．モゲという概念を用いる~\cite{goto2010bioruby}．
モゲという概念を用いる．モゲという概念を用いる．モゲという概念を用いる．
モゲという概念を用いる~\cite{matsumoto2002ruby}．モゲという概念を用いる．モゲという概念を用いる．
モゲという概念を用いる．モゲという概念を用いる．モゲという概念を用いる．
モゲという概念を用いる．モゲという概念を用いる~\cite{fulton2006ruby}．

モゲという概念を用いる．モゲという概念を用いる．モゲという概念を用いる．
モゲという概念を用いる．モゲという概念を用いる．モゲという概念を用いる．
モゲという概念を用いる．モゲという概念を用いる．モゲという概念を用いる．
モゲという概念を用いる．モゲという概念を用いる．モゲという概念を用いる．
図~\ref{fig:rectangle}の長方形に角は四つである。
本文中から必ず参照する。オプションについても調べた。
モゲという概念を用いる．モゲという概念を用いる．モゲという概念を用いる．
モゲという概念を用いる．モゲという概念を用いる．モゲという概念を用いる．

PowerPoint で図を作るときは，スライドサイズを先に指定すると余白が少なくなる．

% setsumei
\begin{figure}[htbp]
\centering
\includegraphics[width=10cm]{img/sample.pdf}
\caption{PowerPoint 上での矩形}
\label{fig:rectangle}
\end{figure}%

モゲという概念を用いる．モゲという概念を用いる．モゲという概念を用いる．
モゲという概念を用いる．モゲという概念を用いる．モゲという概念を用いる．
モゲという概念を用いる．モゲという概念を用いる．モゲという概念を用いる．

\section{システム設計}
\label{subsec:app_yy}

模気という構造を用いる．模気という構造を用いる．模気という構造を用いる~\cite{richardson2008restful}．
模気という構造を用いる．模気という構造を用いる．模気という構造を用いる．
模気という構造を用いる．模気という構造を用いる．模気という構造を用いる．
表~\ref{tab:情テク}は豹ではなく表である。
本文中から必ず参照する。オプションについては解説しない。

% setsumei
\begin{table}[htbp]
 \caption{情報テクノロジー学科線路側研究室○○情報}
 \label{tab:情テク}
 \centering
  \begin{tabular}{l|lr}
  \hline
  研究室 & 場所 & 人数\\
  \hline
  佐久田研 & O-YYY & 5 \\
  原田研 & O-XXX & 7 \\
  D\"urst 研 & O-527 & 3\\
  大原研 & O-ZZZ & 2 \\
  \hline
 \end{tabular}
\end{table}

模気という構造を用いる．模気という構造を用いる．模気という構造を用いる．
模気という構造を用いる．模気という構造を用いる．模気という構造を用いる．
模気という構造を用いる．模気という構造を用いる．模気という構造を用いる~\cite{sasada2005yarv}．
模気という構造を用いる．模気という構造を用いる．模気という構造を用いる．
模気という構造を用いる．模気という構造を用いる．模気という構造を用いる．
模気という構造を用いる．模気という構造を用いる．模気という構造を用いる．
模気という構造を用いる．模気という構造を用いる．模気という構造を用いる．
模気という構造を用いる．模気という構造を用いる．模気という構造を用いる．
